\SourceLeaf{125}
\fancyhead[L]{\small Allgemeine Theorie des Erdmagnetismus}
\renewcommand{\Article}[1]{\par\medskip\Needspace{4\baselineskip}\pdfbookmark[1]{Art. #1}{erdmagnetismus-art-#1}\centerline{\textbf{#1.}}\nobreak\smallskip}
\pdfbookmark[0]{Allgemeine Theorie des Erdmagnetismus}{erdmagnetismus-body}
\par\medskip
\centerline{ALLGEMEINE THEORIE}
\centerline{DES}
\centerline{ERDMAGNETISMUS.}
\medskip
Der unermüdliche Eifer, mit dem man in neuerer Zeit in allen Teilen der Erdoberfläche Richtung und Stärke der magnetischen Kraft der Erde zu erforschen sucht, ist umso erfreulicher, je deutlicher dabei das rein wissenschaftliche Interesse hervortritt. So wichtig die möglichst vollständige Kenntnis der Deklinationslinie für die Schifffahrt auch ist, ihr Bedarf reicht nicht weiter, und was darüber hinausgeht, ist für sie beinahe ohne Belang. Die Wissenschaft fördert zwar gern auch materielle Interessen, lässt sich aber nicht darauf beschränken und verlangt für alle Gegenstände ihrer Forschung die gleiche Anstrengung.

Die Ergebnisse magnetischer Beobachtungen stellt man auf Erdkarten gewöhnlich durch drei Liniensysteme dar, die man als Isogonen, Isoklinen und Isodynamen bezeichnet hat. Form und Lage dieser Linien ändern sich im Laufe der Zeit erheblich, sodass eine einzelne Zeichnung den Zustand der Erscheinung nur für einen bestimmten Zeitpunkt angibt. \textsc{Halley}s Deklinationskarte unterscheidet sich stark von \textsc{Barlow}s Darstellung aus dem Jahr 1833, und \textsc{Hansteen}s Inklinationskarte für 1780 weicht bereits stark von der heutigen Lage der Isoklinen ab. Die Versuche, die Intensität darzustellen, sind noch zu neu, als dass sich bei ihr schon ähnliche Änderungen \SourceLeaf{126}nachweisen ließen. Zweifellos werden sie sich im Laufe der Zeit einstellen. Alle diese Karten sind noch mehr oder weniger lückenhaft oder stellenweise unzuverlässig. Auch wenn sie wegen der Unzugänglichkeit einiger Teile der Erdoberfläche nie ganz vollständig sein können, bleibt zu hoffen, dass sie der Vollständigkeit mit raschen Schritten näherkommen.

Vom höheren Standpunkt der Wissenschaft aus ist diese möglichst vollständige Zusammenstellung der beobachteten Erscheinungen aber noch nicht das eigentliche Ziel. Man hat damit nur Ähnliches getan wie der Astronom, wenn er etwa die scheinbare Bahn eines Kometen auf der Himmelskugel beobachtet hat. Solange man die verwickelten Erscheinungen keinem einheitlichen Prinzip unterstellt hat, besitzt man nur die Bausteine und noch kein Gebäude. Für den Astronomen beginnt die Hauptarbeit erst, nachdem ihm der Himmelskörper aus dem Blick geraten ist. Gestützt auf das Gravitationsgesetz berechnet er aus den Beobachtungen die Elemente der wahren Bahn und kann dadurch sogar den weiteren Lauf mit Sicherheit angeben. Ebenso sollte der Physiker sich die Aufgabe stellen, die Grundkräfte, die die Erscheinungen des Erdmagnetismus hervorbringen, nach ihrer Wirkungsweise und ihren Beträgen zu erforschen, wenigstens soweit es die andersartigen und zum Teil weniger günstigen Umstände erlauben. Er sollte die Beobachtungen, soweit sie reichen, diesen Elementen unterordnen und so die Erscheinungen auch für Gebiete, in die die Beobachtung nicht vordringen konnte, wenigstens mit einer gewissen verlässlichen Annäherung vorwegnehmen. Es ist jedenfalls gut, dieses höchste Ziel vor Augen zu behalten und zu versuchen, die Wege dorthin gangbar zu machen, auch wenn die verfügbaren Daten gegenwärtig so unvollkommen sind, dass man sich dem Ziel selbst nur von Ferne nähern kann.

Ich will hier die früheren erfolglosen Versuche nicht erörtern, bei denen man das große Rätsel ohne jede physikalische Grundlage erraten zu können glaubte. Eine physikalische Grundlage kann man nur solchen Versuchen zugestehen, die die Erde als wirklichen Magneten betrachten und der Rechnung allein die nachgewiesene Fernwirkung eines Magneten zugrunde legen. Alle bisherigen Versuche dieser Art haben jedoch eines gemeinsam. Statt zunächst zu untersuchen, \textit{wie} dieser große Magnet beschaffen sein muss, um den Erscheinungen gerecht zu werden, und sich dabei auf eine \SourceLeaf{127}einfache oder eine sehr komplizierte Beschaffenheit einzustellen, gingen sie von vornherein von einer bestimmten einfachen Beschaffenheit aus und prüften, ob die Erscheinungen mit dieser Hypothese vereinbar waren. Darin wiederholt sich indessen nur, was die Geschichte der Astronomie und der Naturwissenschaften über die Anfänge so vieler unserer Kenntnisse berichtet.

Die einfachste Hypothese dieser Art nimmt einen einzigen sehr kleinen Magneten im Mittelpunkt der Erde an. Genauer gesagt setzt sie, da wohl kaum jemand ernsthaft an die wirkliche Existenz eines solchen Magneten geglaubt hat, voraus, dass der Magnetismus in der Erde so verteilt ist, dass seine Gesamtwirkung nach außen der Wirkung eines gedachten unendlich kleinen Magneten entspricht. Das lässt sich etwa mit der Gravitation zu einer homogenen Kugel vergleichen. Sie entspricht der Anziehung durch eine gleich große, im Mittelpunkt konzentrierte Masse. Unter dieser Voraussetzung sind die beiden Punkte, in denen die verlängerte magnetische Achse dieses Zentralmagneten die Erdoberfläche schneidet, die magnetischen Pole der Erde. Dort steht die Magnetnadel senkrecht, und zugleich ist die Intensität am größten. Auf dem Großkreis in der Mitte zwischen beiden Polen, dem magnetischen Äquator, ist die Inklination $=0$ und die Intensität halb so groß wie an den Polen. Zwischen dem magnetischen Äquator und einem Pol hängen sowohl Inklination als auch Intensität nur vom Abstand von diesem Äquator ab, also von der magnetischen Breite. Dabei ist der Tangens der Inklination gleich dem Doppelten des Tangens dieser Breite. Schließlich fällt die Richtung der horizontalen Nadel überall mit der Richtung eines zum nördlichen magnetischen Pol gezogenen Großkreises zusammen. Mit all diesen notwendigen Folgen der Hypothese stimmt die Natur jedoch nur ungefähr überein. In Wirklichkeit ist die Linie verschwindender Inklination eine doppelt gekrümmte Linie und kein Großkreis. Gleiche Neigungen ergeben keine gleichen Intensitäten. Die Richtungen der horizontalen Nadel laufen längst nicht alle in einem einzigen Punkt zusammen, und so weiter. Schon eine oberflächliche Betrachtung genügt also, um zu zeigen, dass diese Hypothese zu verwerfen ist. Trotzdem verwendet man einen der obigen Sätze noch heute als Näherung, um die Lage der Linie verschwindender Inklination aus Beobachtungen abzuleiten, die in einiger Entfernung von ihr bei mäßigen Inklinationen gemacht wurden.

Von einer ähnlichen Hypothese war bereits vor 80 Jahren \textsc{Tobias Mayer} \SourceLeaf{128}ausgegangen. Er setzte den unendlich kleinen Magneten allerdings in einem Abstand von etwa einem Siebtel des Erdradius vom Mittelpunkt an. Vermutlich wollte er die Rechnung nicht weiter verkomplizieren und behielt daher die an sich völlig willkürliche Einschränkung bei, dass die zur Magnetachse senkrechte Ebene durch den Erdmittelpunkt geht. So fand er, freilich nur für eine sehr kleine Zahl von Orten, eine recht gute Übereinstimmung der beobachteten Deklinationen und Inklinationen mit seiner Rechnung. Eine umfassendere Prüfung hätte aber bald gezeigt, dass sich mit dieser Hypothese die Gesamtheit der Erscheinungen bei diesen beiden Elementen nicht viel besser darstellen lässt als mit der zuerst erwähnten. Intensitätsbestimmungen gab es damals bekanntlich noch gar nicht.

\expandafter\def\csname HistoricalNoteText14\endcsname{Bei der Deklination kommt sogar einmal eine Abweichung von mehr als 29 Grad vor. Gerechterweise sollte man den Rechenfehler aber nach dem tatsächlichen Unterschied zwischen der berechneten und der beobachteten Gesamtrichtung beurteilen. Die Zahl der Deklinationsgrade allein gibt diesen Unterschied nicht an. Am betreffenden Ort beträgt er $11\frac{1}{2}$ Grad.}
\textsc{Hansteen} ist einen Schritt weiter gegangen und hat versucht, die Hypothese \textit{zweier} unendlich kleiner Magnete von ungleicher Lage und Stärke den Erscheinungen anzupassen. Entscheidend für die Zulässigkeit einer Hypothese bleibt immer der Vergleich ihrer Ergebnisse mit der Erfahrung. \textsc{Hansteen} hat seine Hypothese mit Beobachtungen an 48 verschiedenen Orten verglichen, darunter jedoch nur 12, an denen auch die Intensität bestimmt wurde, und überhaupt nur 6, für die alle drei Elemente vorliegen. Wir finden hier noch Abweichungen zwischen Rechnung und Beobachtung, die bei der Inklination auf fast 13 Grad ansteigen\HistoricalNote{14}.

Hält man so große Abweichungen für unvereinbar mit den Anforderungen an eine zufriedenstellende Theorie, so muss man folgern, dass die magnetische Beschaffenheit des Erdkörpers nicht durch eine Konzentration auf einen oder zwei einzelne unendlich kleine Magnete dargestellt werden kann. Damit ist nicht ausgeschlossen, dass man mit einer \textit{größeren} Anzahl solcher gedachten Magnete schließlich eine genügende Übereinstimmung erreichen könnte. Ob eine solche Form der Lösung ratsam wäre, ist aber eine ganz andere Frage. Tatsächlich scheint es, dass die schon bei zwei Magneten so \SourceLeaf{129}überaus mühsamen Rechnungen bei einer erheblich größeren Zahl auf unüberwindliche praktische Schwierigkeiten stoßen würden. Am besten wird es sein, diesen Weg ganz zu verlassen, der unwillkürlich an die Versuche erinnert. die Planetenbewegungen durch immer mehr Epizyklen zu erklären.

In dieser Abhandlung werde ich die allgemeine Theorie des Erdmagnetismus unabhängig von allen besonderen Hypothesen über die Verteilung der magnetischen Flüssigkeit im Erdkörper entwickeln und zugleich die Ergebnisse mitteilen, die ich bei der ersten Anwendung der Methode erhalten habe. So unvollkommen diese Ergebnisse auch sein müssen, können sie doch eine Vorstellung davon geben, was man künftig zu erreichen hoffen darf, wenn erst zuverlässige und vollständige Beobachtungen aus allen Teilen der Erde zur Verfügung stehen, um die Ergebnisse genauer und wiederholt zu verfeinern.

